% ============================================
% Appendix
% Outside 8-page limit
% ============================================

\section{Evaluation archive and aggregation rule}
\label{app:configs}

The archive is organized so that each reported win rate can be
recomputed from completed-run records. It includes one record per
trajectory, the condition tags used to select cells, frozen $L_4{+}L_5$
snapshots for the within-codebase ablation, representative
decision-time and postrun prompt records, and scripts for Wilson 95\%
win-rate intervals and bootstrap 95\% score intervals. A cell's win
rate is simply victories divided by completed games in that cell. We
keep the fixed-$A_0$, cross-backbone, ladder, and full-archive
streams separate. Table~\ref{tab:evidence-streams} gives the
denominator rule for each stream.

\begin{table}[H]
\centering\footnotesize
\caption{Denominator map. Only the balanced fixed-$A_0$ subset enters the headline ablation (Table~\ref{tab:fivecond}); streams are never pooled.}
\label{tab:evidence-streams}
\setlength{\tabcolsep}{4pt}
\renewcommand{\arraystretch}{1.08}
\begin{tabular*}{\columnwidth}{@{\extracolsep{\fill}}llll@{}}
\toprule
Stream & $N$ & Metric & Role \\
\midrule
Fixed-$A_0$ & 50 $(5\!\times\!10)$ & win+score & headline \\
Backbone & 5/cell & score shift & diagnostic \\
Ladder & unequal & max $A$ & endpoint \\
Archive & 298 & tags+scripts & audit \\
\bottomrule
\end{tabular*}
\end{table}
\FloatBarrier

\paragraph{Score-formula audit.}
The derived score (Eq.~\ref{eq:score}) uses $\mathrm{bosses}=0$
when floor $<18$, $1$ when floor $<34$, $2$ otherwise, and $3$
for victories. Table~\ref{tab:cell-floor-boss} lets readers
reproduce the paper-reported means.

\paragraph{Confidence intervals.}
For a cell with $w$ wins out of $n$ completed runs, the Wilson 95\%
interval \citep{wilson1927} on the win rate $\hat p = w/n$ is
\begin{equation}
\frac{\hat p + \tfrac{z^2}{2n} \pm z\sqrt{\tfrac{\hat p(1-\hat p)}{n} + \tfrac{z^2}{4n^2}}}{1 + z^2/n},
\quad z=1.96.
\label{eq:wilson}
\end{equation}
Score intervals use the percentile bootstrap \citep{efron1979}: 5{,}000
resamples with replacement from the cell's $n$ run-level scores,
reporting the empirical $[2.5,97.5]$ percentiles. The descriptive
pooled scaffolded row uses an exact Clopper--Pearson interval and is
labeled as such.

\begin{table}[H]
\centering\footnotesize
\caption{Per-cell floor and boss-clear counts that reproduce the
mean scores reported in Table~\ref{tab:fivecond}.}
\label{tab:cell-floor-boss}
\setlength{\tabcolsep}{4pt}
\renewcommand{\arraystretch}{1.05}
\begin{tabular*}{\columnwidth}{@{\extracolsep{\fill}}lccccc@{}}
\toprule
Cell & $N$ & Wins & $\overline{\mathrm{floor}}$ & $\overline{\mathrm{bosses}}$ & $\overline{\mathrm{score}}$ \\
\midrule
baseline-strict & 10 & 3 & 39.2 & 1.80 & 70.40 \\
prompt-only     & 10 & 4 & 38.4 & 1.80 & 69.60 \\
mode-a          & 10 & 6 & 43.9 & 2.40 & 85.50 \\
mode-b-frozen   & 10 & 6 & 43.4 & 2.30 & 83.27 \\
full-frozen     & 10 & 6 & 42.2 & 2.30 & 82.07 \\
\bottomrule
\end{tabular*}
\end{table}
\FloatBarrier

\section{Memory contract: five typed layers}
\label{app:memory-contract}

Table~\ref{tab:memory-contract} gives a compact legend for the five
stores used in the paper. The useful distinction is mutability:
$L_1$ and $L_2$ are fixed, $L_3$ can be filtered, and $L_4/L_5$ can be
disabled, frozen, or made writable depending on the evidence stream.

\begin{table}[H]
\centering\footnotesize
\caption{Five-layer memory contract. Each decision receives typed slices; raw cross-decision transcript is not appended.}
\label{tab:memory-contract}
\setlength{\tabcolsep}{4pt}
\renewcommand{\arraystretch}{1.08}
\begin{tabular*}{\columnwidth}{@{\extracolsep{\fill}}cllll@{}}
\toprule
Layer & Store & Key & Write & Ablation \\
\midrule
$L_1$ & protocol & state & fixed & always \\
$L_2$ & schema & decision & fixed & always \\
$L_3$ & rules & entities & refresh & filter \\
$L_4$ & episodes & char/$A$/act & postrun & off/on \\
$L_5$ & skills & trigger & gated & off/A/B \\
\bottomrule
\end{tabular*}
\end{table}
\FloatBarrier

\section{Concrete prompt and learning examples}
\label{app:artifacts}

\paragraph{Decision-time composed prompt.}
At decision time, the agent assembles the user prompt in a fixed
typed order: fired $L_5$ skills, $L_4$ episodes, $L_3$ game facts,
the $L_2$ typed state prompt, and a schema hint listing valid
actions. The order is shared across non-combat decision states. Combat
is the only intra-fight stateful container, and its emitted message
list per round is bounded to three items: combat start, an
acknowledgement, and the latest user state. Earlier rounds are not
copied into the transcript; the next round is prompted from a newly
composed user message. A shortened floor-7 card-reward excerpt shows
the typed layers:

\begin{Verbatim}[fontsize=\footnotesize]
## Expert Knowledge (retrieved skills)
**Silent - Draft and Shop Rules** (seed)
Early rewards must solve damage first,
then block, draw, and energy. ...

## Card-Specific Insights
- dodge and roll: delayed block plus Dexterity.
- slice: 0-cost transitional damage.

## Game Knowledge / Card Mechanics
- DodgeAndRoll: Block Next Turn; upgrade +2.
- Slice: Upgrade: Damage +3

## Card Reward
HP: 57/57 (100%) | Gold: 55 | Act: 1 | Floor: 7

## Decision Format (card_reward_action)
Valid actions: choose_reward_card
or choose_reward_alternative
\end{Verbatim}

\paragraph{Representative learned skill.}
One learned skill concerns a boss mechanic that exhausts attacks and
skills on specific turns (turns 2, 5, 8, and 11). The stored rule
instructs the agent to preserve core scaling cards on those turns and
to prefer powers or natively exhausting cards. The example shows that
$L_5$ stores state-conditioned tactical rules with Boolean triggers,
not generic advice or similarity-retrieved raw logs. Full prompt and
learned-skill records are included in the artifact archive.

\section{External calibration and source boundaries}
\label{app:difficulty-calibration}

External \stspad{} statistics enter the paper as difficulty and
ecosystem context. Cached snapshots of each source are released
with the artifact archive so that follow-up work can reproduce
the calibration rows under identical inputs.

\begin{table}[H]
\centering\footnotesize
\caption{External numbers calibrate difficulty and ecosystem context only.}
\label{tab:difficulty-calibration}
\setlength{\tabcolsep}{4pt}
\renewcommand{\arraystretch}{1.08}
\begin{tabular*}{\columnwidth}{@{\extracolsep{\fill}}lll@{}}
\toprule
Source & Role & Not used as \\
\midrule
AGI-Eval & LLM yardstick & causal baseline \\
Mega Crit $A_0$ & difficulty anchor & matched human test \\
Spiracle & community context & population rate \\
\bottomrule
\end{tabular*}
\end{table}
\FloatBarrier

% The new "Full decision-time prompt exhibits" section (appendix_two_full_prompt_exhibits.tex)
% is \input{} after this file from main_v2.tex and provides complete prompt
% content for both a combat and a shop decision.
